\section{Evaluation}
\label{sec:evaluation}

The introduction claimed five contributions.
This section evaluates each against the benchmark evidence.

\textbf{(i) Unified pipeline.}
All seven benchmark scenarios (Table~\ref{tab:scenarios}) ran from a single YAML configuration through one binary.
No external tool or intermediate file format was required.
The claim holds.

\textbf{(ii) First genome-wide UMI/duplex simulator.}
The panel UMI scenario (1\,MB, 200x, simplex) completed successfully (Table~\ref{tab:scenarios}).
hg38 chr22 configs for UMI simplex and Twist-style duplex (both at 30x) are included in the hg38 benchmark suite and complete without error.
No other published tool generates genome-wide UMI-tagged reads with configurable family sizes and duplex strand pairing.
The claim holds.

\textbf{(iii) First cfDNA generator with variant support.}
The cfDNA scenario (200x, 2\% purity, cfDNA fragment model) produced reads with nucleosomal fragment distributions and injected somatic variants.
No other tool combines these two features.
The claim holds.

\textbf{(iv) Stochastic VAF.}
All variant-containing benchmarks used per-read Bernoulli sampling.
Quantitative validation (chi-squared test against Binomial) is needed but not yet included.
The claim is supported by design but not yet independently validated.

\textbf{(v) Streaming Rust with bounded memory.}
Per-pair memory is constant at ${\sim}0.45$\,KB across coverage depths on the 1\,MB reference (Figure~\ref{fig:mem}).
The streaming architecture bounds memory by concurrent batch state rather than total dataset size.
On chr22 at 30x (Table~\ref{tab:hg38}), per-pair memory rises to ${\sim}0.85$\,KB because the larger chromosome produces more pairs per chunk.
The claim holds in the sense that memory is bounded by channel depth, not dataset size.
The bound itself grows with chromosome length and coverage; the \texttt{output\_buffer\_regions} setting provides a lever to reduce it on memory-constrained systems.

\textbf{Thread scaling limitation.}
Adding threads degrades wall time for the 10\,MB / 30x I/O-bound workload (1 thread = 140\,s; 12 threads = 160\,s).
This is an honest result: FASTQ compression saturates the single writer thread, and coordination overhead exceeds compute savings.
Compute-bound workloads (UMI simulation, high feature density) would scale better, but this was not benchmarked across thread counts.
